% =========================================================
% Computational Linear Algebra — The Matrix of a Linear Transformation
% (Strang §8.2 / Johnston §1.4 / Ashraf §3.5–3.6)
% =========================================================

\section{The Matrix of a Linear Transformation}

\begin{tcolorbox}[colback=gray!6, colframe=gray!40, arc=2pt,
  left=6pt, right=6pt, top=4pt, bottom=4pt,
  title={\small\textbf{Standard Matrix --- Quick Reference}},
  fonttitle=\small\bfseries]
\begin{tabular}{@{}p{0.30\textwidth}p{0.63\textwidth}@{}}
\textbf{Standard matrix}      & $[T] = [T(\mathbf{e}_1)\mid T(\mathbf{e}_2)\mid\cdots\mid T(\mathbf{e}_n)] \in \mathbb{R}^{m\times n}$. \\[4pt]
\textbf{How to build it}       & Evaluate $T$ on each standard basis vector; place results as columns. \\[4pt]
\textbf{Uniqueness}            & $[T]$ is the \emph{unique} matrix with $T(\mathbf{v}) = [T]\mathbf{v}$ for all $\mathbf{v}$. \\[4pt]
\textbf{Composition}           & $[S \circ T] = [S][T]$. \\[4pt]
\textbf{Change of basis}       & Same $T$, different bases $\Rightarrow$ different matrix; related by $B = P^{-1}AP$. \\
\end{tabular}
\end{tcolorbox}

% ---------------------------------------------------------
\subsection{The Standard Matrix}

\begin{tcolorbox}[colback=thmbox, colframe=thmborder, arc=2pt,
  left=6pt, right=6pt, top=4pt, bottom=4pt,
  title={\small\textbf{Theorem (Standard Matrix of a Linear Transformation)}},
  fonttitle=\small\bfseries]
A function $T : \mathbb{R}^n \to \mathbb{R}^m$ is a linear transformation if and only if
there exists a unique matrix $[T] \in \mathbb{R}^{m \times n}$ such that
\[
  T(\mathbf{v}) = [T]\mathbf{v} \quad \text{for all } \mathbf{v} \in \mathbb{R}^n.
\]
The \emph{standard matrix} $[T]$ is constructed by placing the images of
the standard basis vectors as columns:
\[
  [T] = \bigl[\, T(\mathbf{e}_1) \;\mid\; T(\mathbf{e}_2) \;\mid\; \cdots \;\mid\; T(\mathbf{e}_n) \,\bigr].
\]
\end{tcolorbox}

\begin{remark*}[Why this works]
Since every $\mathbf{v} = v_1\mathbf{e}_1 + \cdots + v_n\mathbf{e}_n$,
linearity gives
\[
  T(\mathbf{v}) = v_1 T(\mathbf{e}_1) + \cdots + v_n T(\mathbf{e}_n) = [T]\mathbf{v}.
\]
Knowing $T$ on the basis vectors determines $T$ on the entire space.
\end{remark*}

% ---------------------------------------------------------
\subsection{Matrix Under Change of Basis}

\begin{tcolorbox}[colback=propbox, colframe=propborder, arc=2pt,
  left=6pt, right=6pt, top=4pt, bottom=4pt,
  title={\small\textbf{Definition (Matrix Relative to a Basis)}},
  fonttitle=\small\bfseries]
Let $T : \mathbb{R}^n \to \mathbb{R}^n$ be a linear transformation and let
$\mathcal{B} = \{\mathbf{b}_1, \ldots, \mathbf{b}_n\}$ be an ordered basis.
The \emph{matrix of $T$ relative to $\mathcal{B}$}, denoted $[T]_\mathcal{B}$,
has as its $j$-th column the coordinates of $T(\mathbf{b}_j)$ expressed in the
basis $\mathcal{B}$.

If $P$ is the transition matrix whose columns are $\mathbf{b}_1, \ldots, \mathbf{b}_n$
expressed in the standard basis, then
\[
  [T]_\mathcal{B} = P^{-1} [T] P.
\]
Two matrices related by $B = P^{-1}AP$ are called \emph{similar}.
\end{tcolorbox}

\begin{remark*}[NURBS connection: Bézier curve in matrix form]
The Bézier curve has the matrix representation
\[
  C(u) = [u^i]^T M [P_i],
\]
where:
\begin{itemize}
  \item $[u^i]^T = (1, u, u^2, \ldots, u^n)$ is the monomial basis vector,
  \item $[P_i]$ is the column vector of control points,
  \item $M$ is the $(n+1)\times(n+1)$ matrix that encodes the Bernstein polynomials
        in the monomial basis.
\end{itemize}
The matrix $M$ is the standard matrix of the linear transformation
that converts from Bernstein basis coordinates to monomial basis coordinates.
Setting $[\mathbf{a}_i] = M[P_i]$ gives the power basis representation of the
same curve (so the Bézier coefficients $[P_i]$ are the coordinates of the
curve in the Bernstein basis, and $[\mathbf{a}_i]$ are the coordinates in the
monomial basis).
Converting between the two is exactly the change-of-basis operation $P^{-1}AP$
applied to the coefficient vectors.
\end{remark*}
